\newcommand{\nr}{N}
\newcommand{\set}[1]{\left\{ \def\given{\ \middle| \ }  #1 \right\}  }
\newcommand{\defeq}{\mathrel{:=}}
\NewDocumentCommand{\iso}{}{\cong}
\DeclareMathOperator{\tr}{tr}
\DeclareMathOperator{\id}{id}
\DeclareMathOperator{\End}{End}
\renewcommand{\hom}{\operatorname{Hom}}

\NewDocumentCommand{\nset}{m}{\left[#1\right]}

\renewcommand{\Re}{\operatorname{Re}}
\renewcommand{\Im}{\operatorname{Im}}


\NewDocumentCommand{\dif}{}{\,\mathrm{d}}

\NewDocumentCommand{\bigoh}{}{\mathrm{O}}

\NewDocumentCommand{\du}{}{\amalg}

%use ifempty from xparse
\ExplSyntaxOn
\DeclareExpandableDocumentCommand{\IfEmptyTF}{mmm}
 {
    \tl_if_empty:nTF {#1} {#2} {#3}
 }
\ExplSyntaxOff

%fix spacing with paried delimiters
%use these for function arguments
\def\leftfun#1{\mathopen{}\left#1}
\def\rightfun#1{\right#1}

%useful for defining named functions
\NewDocumentCommand{\ParenIfNonempty}{m}{%
  \IfEmptyTF{#1}{%
  }{%
    \leftfun(#1\rightfun)%
  }
}

\NewDocumentCommand{\ii}{}{\mathsf{i}}
\NewDocumentCommand{\tu}{}{2 \pi \ii}

\NewDocumentCommand{\evec}{m}{\mathbf{e}_{#1}}

\NewDocumentCommand{\slg}{}{\operatorname{SL}_2(\mathbb{C})}
\NewDocumentCommand{\sla}{}{\mathfrak{sl}_2}
\NewDocumentCommand{\pslg}{}{\operatorname{PSL}_2(\mathbb{C})}
\NewDocumentCommand{\pjsp}{O{\mathbb{C}}}{\operatorname{#1 P}^{1}}
\DeclareMathOperator{\hyp}{hyp}
\NewDocumentCommand{\qn}{}{\triangleleft}
\NewDocumentCommand{\qninv}{}{\mathbin{\overline{\triangleleft}}}

\NewDocumentCommand{\qgrplong}{O{\xi}}{\mathcal{U}_{#1}(\mathfrak{sl}_2)}
\NewDocumentCommand{\qgrp}{O{\xi}}{\mathcal{U}_{#1}}
\NewDocumentCommand{\weyl}{O{\xi}}{\mathcal{W}_{#1}}

\NewDocumentCommand{\mer}{}{\mathbf{m}}
\NewDocumentCommand{\lon}{}{\mathbf{l}}
\NewDocumentCommand{\extr}{m}{M_{#1}}
\NewDocumentCommand{\comp}{m}{S^3 \setminus #1}
\DeclareMathOperator{\lk}{Lk}

\NewDocumentCommand{\homol}{m m}{\operatorname{H}_{#1}\leftfun(#2\rightfun)}
\NewDocumentCommand{\cohomol}{m m}{\operatorname{H}^{#1}\leftfun(#2\rightfun)}

%all Reidemeister moves in one macro
\ExplSyntaxOn
\NewDocumentCommand{\reid}{m O{}}{
  \bool_case:nTF
  {
    {\str_if_eq_p:NN {#1} {1} } { \operatorname{R}\sb{1} }
    {\str_if_eq_p:NN {#1} {1f} } { \operatorname{R}\sb{1}^{\operatorname{fr}} }
    {\str_if_eq_p:NN {#1} {2} } { \operatorname{R}\sb{2}^{#2} }
    {\str_if_eq_p:NN {#1} {2a} } { \operatorname{R}\sb{2}^{\operatorname{a}} }
    {\str_if_eq_p:NN {#1} {2b} } { \operatorname{R}\sb{2}^{\operatorname{b}} }
    {\str_if_eq_p:NN {#1} {3} } { \operatorname{R}\sb{3} }
  }
  { } %matches
  { } %unmatched
}
\ExplSyntaxOff

\NewDocumentCommand{\modcat}{m}{#1\operatorname{-mod}}
\NewDocumentCommand{\tsunit}{}{\mathbf{1}}

\NewDocumentCommand{\sclbrak}{m}{\left\langle #1 \right \rangle}

\NewDocumentCommand{\point}{}{\cdot}

\NewDocumentCommand{\repvar}{m}{\mathsf{R}\leftfun(#1\rightfun)}
\NewDocumentCommand{\admvar}{m}{\mathsf{A}\leftfun(#1\rightfun)}
\NewDocumentCommand{\logvar}{m}{\mathfrak{F}\leftfun(#1\rightfun)}
\NewDocumentCommand{\admvarbdry}{m}{\mathsf{A}^{\partial}_{#1}}

\NewDocumentCommand{\csbun}{O{}}{\mathcal{E}^{#1}}
\NewDocumentCommand{\qbun}{O{\nr} O{}}{\mathcal{E}_{#1}^{#2}}
\NewDocumentCommand{\qbunbig}{O{\nr} O{}}{\tilde{\mathcal{E}}_{#1}^{#2}}

\NewDocumentCommand{\qfunc}{O{\nr} m}{\mathcal{Z}_{#1}^{\psi}\ParenIfNonempty{#2}}

\NewDocumentCommand{\qinv}{O{\nr} m}{\mathrm{Z}_{#1}^{\psi}\ParenIfNonempty{#2}}

\NewDocumentCommand{\actq}{O{\nr} m m}{%
  \mathcal{S}_{#1}\leftfun(#2
  \IfEmptyTF{#3}{%
    \rightfun)
  }{%
    \middle|#3\rightfun)%
  }
}

\NewDocumentCommand{\act}{ m m}{%
  \mathcal{S}\leftfun(#1
  \IfEmptyTF{#2}{%
    \rightfun)
  }{%
    \middle|#2\rightfun)%
  }
}


\NewDocumentCommand{\stint}{O{\nr} m m}{%
  J_{#1}\leftfun(#2
  \IfEmptyTF{#3}{%
    \rightfun)
  }{%
    \middle|#3\rightfun)%
  }
}

\NewDocumentCommand{\stintadj}{O{\nr} m m}{%
  \widetilde{J}_{#1}\leftfun(#2
  \IfEmptyTF{#3}{%
    \rightfun)
  }{%
    \middle|#3\rightfun)%
  }
}

\NewDocumentCommand{\crossingfunction}{O{\nr} m m m}{
  \Omega_{#1}^{#2}\leftfun(#3
  \IfEmptyTF{#4}{%
    \rightfun)
  }{%
    \middle|#4\rightfun)%
  }
}

\NewDocumentCommand{\crossingfunctionclass}{m m}{
  \Omega^{#1}
  \IfEmptyTF{#2}{%
  }{%
    \leftfun(#2\rightfun)%
  }
}

\NewDocumentCommand{\lkfunc}{O{\nr} m}{\mathcal{L}_{#1}\ParenIfNonempty{#2}}
\NewDocumentCommand{\lklog}{O{} m}{\psi_{#1}\ParenIfNonempty{#2}}

\NewDocumentCommand{\qcat}{O{\nr}}{\mathsf{C}_{#1}}
\NewDocumentCommand{\tcat}{O{\nr}}{\mathsf{T}_{#1}}
\NewDocumentCommand{\vect}{}{\mathsf{Vect}}

\NewDocumentCommand{\shadset}{}{\mathsf{U}}

% (quantum) dilogs and q-series stuff
\NewDocumentCommand{\pf}{O{} m }{\varphi_{#1}\ParenIfNonempty{#2}}
\NewDocumentCommand{\pfl}{O{\nr} m}{\phi_{#1}\ParenIfNonempty{#2}}

\NewDocumentCommand{\qlfl}{O{\nr} m m}{%
  \lambda_{#1}\leftfun(#2
  \IfEmptyTF{#3}{%
    \rightfun)
  }{%
    \middle|#3\rightfun)%
  }
}

\NewDocumentCommand{\dll}{m }{\operatorname{L}\ParenIfNonempty{#1}}

\newcommand{\evup}[1]{\operatorname{ev}^\uparrow_{#1}}
\newcommand{\coevup}[1]{\operatorname{coev}^\uparrow_{#1}}
\newcommand{\evdown}[1]{\operatorname{ev}^\downarrow_{#1}}
\newcommand{\coevdown}[1]{\operatorname{coev}^\downarrow_{#1}}

\NewDocumentCommand{\lN}{}{\mathrm{N}}
\NewDocumentCommand{\lW}{}{\mathrm{W}}
\NewDocumentCommand{\lS}{}{\mathrm{S}}
\NewDocumentCommand{\lE}{}{\mathrm{E}}
\NewDocumentCommand{\up}{}{\uparrow}
\NewDocumentCommand{\dn}{}{\downarrow}

% macros borrowed from Colleen Delaney
\newcommand{\parcen}[1]{\phantom{(}#1\phantom{)}}
\newcommand{\br}{to [out=00,in=180]}
